%----------------------------------------------------------------------------------------
%	PAQUETES Y OTRAS CONFIGURACIONES
%----------------------------------------------------------------------------------------

%----------------------------------------------------------------------------------------
%	PAQUETES Y OTRAS CONFIGURACIONES
%----------------------------------------------------------------------------------------

\documentclass[paper=letter, fontsize=11pt]{scrartcl} % Tamaño de papel y letra para el documento

\usepackage[utf8]{inputenc} % Los caracteres acentuados se pueden escribir normalmente en el código
\usepackage[T1]{fontenc} % Configuración de fuente de salida
\usepackage{fourier} % Se usa una fuente diferente al default
\usepackage[spanish,es-noquoting]{babel} % Se configura como documento en español
\usepackage{amsmath,amsfonts,amsthm} % Paquetes para escribir formulas matemáticas
\usepackage{graphicx} % Paquetes para incluir imágenes

\usepackage{circuitikz}
\usepackage{tikz}
\usetikzlibrary{arrows}

\usepackage{sectsty} % Paquete para configuración de secciones
\allsectionsfont{\centering \normalfont \scshape} % Los títulos de las secciones son centrados, con la misma fuente y pequeñas mayúsculas

\usepackage{todonotes}
\usepackage{microtype}

\usepackage{fancyhdr} % Paquete para personalizar pies y cabeceras de página
\pagestyle{fancyplain} % Todas las páginas con las mismas cabeceras y pies de página
\fancyhead{} % Sin cabecera
\fancyfoot[L]{} % Vacío en la izquierda del pie de página
\fancyfoot[C]{} % Vacío en el centro del pie de página
\fancyfoot[R]{\thepage} % Número de página en el pie de pagina
\renewcommand{\headrulewidth}{0pt} % Sin lineas en la cabecera
\renewcommand{\footrulewidth}{0pt} % Sin lineas en el pie de página
\setlength{\headheight}{13.6pt} % Altura de cabecera

\numberwithin{equation}{section} % Numera ecuaciones en cada sección
\numberwithin{figure}{section} % Numera figuras en cada sección
\numberwithin{table}{section} % Numera tablas en cada sección

\setlength\parindent{0pt} % Quita la indentación de los párrafos

\newcommand{\horrule}[1]{\rule{\linewidth}{#1}} % Comando personalizado para hacer linea horizontal


%----------------------------------------------------------------------------------------
%	TITULO
%----------------------------------------------------------------------------------------

\title{	
	\normalfont \normalsize
	\begin{figure}[h]
		\begin{center}
			\includegraphics[width=0.3\textwidth]{../images/UNITEC.png} % La imagen esta incluida en el mismo directorio del código
		\end{center}
	\end{figure}
	\textsc{Instrumentación y Control} \\ [25pt]
	\horrule{0.5pt} \\[0.4cm] % Linea horizontal delgada
	\huge Práctica 2 - Medición de variables físicas \\ % Titulo de la práctica
	\horrule{2pt} \\[0.5cm] % Linea horizontal mas gruesa
}

\author{Roberto Cadena Vega} % Nombre del profesor

\date{\normalsize 5 de junio de 2014} % Fecha de la práctica

%----------------------------------------------------------------------------------------
%	EMPIEZA EL DOCUMENTO
%----------------------------------------------------------------------------------------

\begin{document}

\maketitle % Imprime el título

%----------------------------------------------------------------------------------------
%	OBJETIVOS
%----------------------------------------------------------------------------------------

\section{Objetivos}

	Implementar un sistema eléctrico simple, que sea capaz de medir la temperatura.

%----------------------------------------------------------------------------------------
%	CONOCIMIENTOS PREVIOS
%----------------------------------------------------------------------------------------

\section{Conocimientos Previos}

%----------------------------------------------------------------------------------------

	\subsection{Sensores}

		Dentro de nuestra materia tenemos distintas tipos de variables que tenemos que controlar para que el proceso se lleve adecuadamente, por lo que necesitamos primero que nada medirlas.

		Los sensores nos ayudarán a hacer esto, pero primero tenemos que aprender a usarlos.

		El sensor que utilizaremos en esta práctica es el LM35, por lo que tendrán que descargar su datasheet, pero tal vez te preguntes que es un datasheet... Bien, pues todos los elementos electrónicos tienen hojas de especificaciones en los que viene toda la información relevante para implementar su uso en un sistema eléctrico.

		Pues lo primero que tienes que hacer es descargar de internet el datasheet de nuestro sensor, una simple búsqueda en Google de "datasheet LM35", te dará como resultado un archivo en formato PDF que puedes descargar e imprimir, para su uso en la práctica.

		Una vez que tengamos nuestro datasheet, lo que tienes que buscar es el voltaje necesario para alimentar nuestro sensor (en caso de que sea necesario). Otro valor importante es la señal que nos regresará el sensor. Otros factores importantes son el tiempo de respuesta, la resistividad interna del elemento, etc., sin embargo, no serán de utilidad en esta práctica.

	\subsection{Acondicionamiento}

		Como te habrás dado cuenta en tu datasheet, el voltaje que nos entrega como salida el LM35 es de $0 V$ a $5 V$ si tomamos en cuenta que el rango de temperaturas que queremos medir es de $0 ^oC$ a $50 ^oC$, pero si quisiéramos meter esto en un sistema industrial, normalmente tendríamos que normalizar este rango a $0 V$ a $24 V$.

		¿Cuál sería el número por el que tenemos que multiplicar, para que la señal del LM35 se convierta a la que queremos? Cuando hablamos de señales que tenemos que acondicionar, a este número le llamamos Ganancia.

		\begin{center}
			\tikzstyle{input} = [coordinate]
			\tikzstyle{output} = [coordinate]
			\tikzstyle{block} = [draw, rectangle, minimum height=3em, minimum width=6em]
			\tikzstyle{init} = [pin edge={to-, thin, black}]

			\begin{tikzpicture}[auto, node distance=3cm, >=latex']
				\node [input, name=entrada] {};
				\node [block, right of=entrada] (sensor) {Señal};
				\node [block, right of=sensor] (amplificador) {Ganancia};
				\node [output, right of=amplificador] (salida) {};

				\draw [->] (entrada) -- node[name=x] {$x$} (sensor);
				\draw [->] (sensor) -- node[name=x] {$u$} (amplificador);
				\draw [->] (amplificador) -- node[name=y] {$y$} (salida);
			\end{tikzpicture}
		\end{center}

		Esta ganancia la vamos a obtener con un sistema eléctrico conocido como amplificador.

		\subsubsection{Amplificador}

			Los amplificadores están compuestos de circuitos integrados llamados Amplificadores Operacionales (u OPAMP's para los cuates), resistencias y cables, cosas que ya sabemos usar.

			El diagrama para un amplificador no inversor es el siguiente:

			\begin{center}
				\begin{circuitikz} \draw
					(0,0) node[op amp](opamp){}
					(opamp.+) to [R, l=$R_I$] (-4, -0.5) to [short, -o] (-4, -0.5) node[left] {$V_I$}
					(opamp.-) to [short, *-] (-2, 0.5) to [R, l=$R_I$] (-5, 0.5) node[ground] {}
					(opamp.-) to [short, *-] (-1.2, 2) to [R, l=$R_f$] (1.2, 2)
					(opamp.out) to [short, *-o] (2, 0) node[right] {$V_O$}
					(opamp.out) to [short, *-] (1.2, 2)
					(opamp.down) node[ground] {}
					(opamp.up) to [short, -o] (-0.08, 1) node[above] {$V_+$}
				;\end{circuitikz}	
			\end{center}
			
			y la ecuación que rige su comportamiento es:

				\begin{equation}
					V_O = \left(\frac{R_f}{R_I} + 1\right) V_I
				\end{equation}

			eso quiere decir que el diagrama queda como sigue:

			\begin{center}
				\tikzstyle{input} = [coordinate]
				\tikzstyle{output} = [coordinate]
				\tikzstyle{block} = [draw, rectangle, minimum height=3em, minimum width=6em]
				\tikzstyle{init} = [pin edge={to-, thin, black}]

				\begin{tikzpicture}[auto, node distance=3cm, >=latex']
					\node [input, name=entrada] {};
					\node [block, right of=entrada] (sensor) {Señal};
					\node [block, right of=sensor] (amplificador) {$ \frac{R_f}{R_I} + 1$};
					\node [output, right of=amplificador] (salida) {};

					\draw [->] (entrada) -- node[name=x] {$x$} (sensor);
					\draw [->] (sensor) -- node[name=x] {$u$} (amplificador);
					\draw [->] (amplificador) -- node[name=y] {$y$} (salida);
				\end{tikzpicture}
			\end{center}

			Al ver esto, notamos que tan solo tenemos que calcular la ganancia, y buscar 2 valores de resistencias que nos dé el valor de ganancia deseado.

			Esto se puede ver de la siguiente manera:

			Recuerda que dependiendo del OPAMP que uses, las terminales de tu circuito integrado serán diferentes, por lo que es absolutamente necesario que tengas el datasheet de tu OPAMP a la mano.

%----------------------------------------------------------------------------------------
%	EQUIPO
%----------------------------------------------------------------------------------------

\section{Equipo}

	El siguiente equipo será proporcionado por el laboratorio, siempre y cuando lleguen en los primeros 15 minutos de la práctica, y hagan el vale conteniendo el siguiente equipo (exceptuando las pinzas).

	\begin{itemize}
		\item 1 Fuente de Alimentación
		\item 1 Multímetro
		\item 1 Cable de alimentación
		\item 2 Cables banana - caimán
		\item Pinzas
	\end{itemize}

%----------------------------------------------------------------------------------------
%	MATERIALES
%----------------------------------------------------------------------------------------

\section{Materiales}

	\begin{itemize}
		\item Protoboard
		\item LM35
		\item LM358
		\item Resistencias
		\begin{itemize}
			\item $180 \Omega$
			\item $220 \Omega$
			\item $330 \Omega$
			\item $1 k\Omega$
			\item $3.3 k\Omega$
			\item $10 k\Omega$
			\item $100 k\Omega$
		\end{itemize}
		\item Cables
	\end{itemize}

%----------------------------------------------------------------------------------------
%	DESARROLLO
%----------------------------------------------------------------------------------------

\section{Desarrollo}
	
	\begin{enumerate}
		\item Diseña un amplificador con la ganancia necesaria para que la salida de nuestro sensor, sea normalizada a un rango de $OV$ a $24V$.
		\item Diseña un amplificador que tenga una ganancia de $-1$ y acopla la salida del amplificador anterior a la entrada de este, para que la señal sea positiva.
		\item Realiza las mediciones requeridas en la hoja de anotaciones.
	\end{enumerate}

%----------------------------------------------------------------------------------------
%	CONCLUSIONES
%----------------------------------------------------------------------------------------

\section{Conclusiones}

	El alumno deberá describir sus conclusiones al final de su reporte de práctica.
    
%----------------------------------------------------------------------------------------
%	HOJA DE ANOTACIONES
%----------------------------------------------------------------------------------------

\clearpage
\section{Hoja de Anotaciones}
	
	\begin{enumerate}
		\item ¿Cuál es la ganancia necesaria para normalizar nuestra señal a la requerida? \newline
		\item ¿Cuál es el voltaje de alimentación que necesita el amplificador? \newline
		\item ¿Cuál es el valor de $R_f$ que necesitas para obtener esa ganancia? \newline
		\item ¿Cuál es el valor de $R_I$ que necesitas para obtener esa ganancia? \newline
		\item ¿Cuál es la temperatura ambiente en el momento de tus mediciones? \newline
		\item ¿Cuál es el valor obtenido al medir con el multímetro, el voltaje de salida en el sensor? \newline
		\item ¿Cuál es el valor obtenido al medir con el multímetro, el voltaje de salida en el primer amplificador? \newline
		\item ¿Cuál es el valor obtenido al medir con el multímetro, el voltaje de salida en el segundo amplificador? \newline
		\item ¿Qué pasa si acercas una fuente de calor al sensor? \newline
	\end{enumerate}

	Integrantes del equipo: \\[0.2cm]
	\horrule{0.5pt} \\[0.2cm] % Linea horizontal delgada
	\horrule{0.5pt} \\[0.2cm] % Linea horizontal delgada
	\horrule{0.5pt} \\[0.2cm] % Linea horizontal delgada
	\horrule{0.5pt} % Linea horizontal delgada	

	Revisó: \\[0.2cm]
	\horrule{0.5pt} \\% Linea horizontal delgada
    
%----------------------------------------------------------------------------------------
%	FIN DEL DOCUMENTO
%----------------------------------------------------------------------------------------

\end{document}